\section{方法选择与结果检验}
\label{sec:ls-direct-linear-method-comparison}

前面各节分别给出了一种求精确解的直接法。面对具体的方阵方程时，要先看矩阵结构选择分解，再用原始数据核验算出的结果。本节按这一顺序归纳。

\paragraph{按结构选择}
非奇异方阵用带部分选主元的LU，对称正定矩阵用Cholesky，对称不定但非奇异的矩阵用分块LDL。三者都先分解、再做三角代入；结构越强，可以省下的运算和存储越多，Cholesky的运算量约为LU的一半，并且不需要选主元。

低秩更新不是新的分解，而是复用上面的求解步骤，并有自己的成立条件：旧矩阵的分解已经存在、修改的秩远小于阶数，并且辅助矩阵$\bm K$可逆且不太病态（第~\ref{sec:ls-woodbury}节）。连续多次更新后，应与重新分解比较代价和精度。

\paragraph{用原始数据核验}
算出$\bm x$后，要代回原方程检查$\bm A\bm x=\bm b$，而不只检查中间的三角方程或辅助方程。例如LU的演示中，$\bm x=(1,-1,1)^\top$代回得$\bm A\bm x=\bm b$；低秩更新的两个例子也都在更新后的矩阵下重算了原残差。浮点计算中残差一般不为零，这时报告缩放残差$\lVert\bm b-\bm A\bm x\rVert/(\lVert\bm A\rVert\lVert\bm x\rVert+\lVert\bm b\rVert)$。

计算中还可以用分解残差（如$\lVert\bm P\bm A-\bm L\bm U\rVert$、$\lVert\bm A-\bm L\bm L^\top\rVert$）核对中间的分解，但这些检查只是旁证，不属于算法本身，也不能代替对原方程的检验。

\paragraph{残差小不等于误差小}
即使上述检验都通过，也不保证$\bm x$的误差小：矩阵接近奇异时，很小的残差可能对应很大的解误差。本章只给出了算法层面的结论：部分选主元控制乘子但不排除主元增长，Cholesky没有增长问题，低秩更新在$\bm K$接近奇异时会放大误差。第~\ref{chap:solution-error-regularization}章用条件数与前向、后向误差系统地解释这些现象（第~\ref{sec:ls-numerical-sensitivity}节），并在直接法与迭代法之间作整体比较（第~\ref{sec:ls-method-selection}节）。
